\section{Related Work}
\label{sec:related}

\paragraph{Cross-vendor GPU benchmarks.}
Extensive literature exists on comparing NVIDIA and AMD GPU performance for
deep learning workloads.  \citet{mlperf_training} and \citet{mlperf_inference}~\citep{mlperf2023}
provide standardized training and inference benchmarks across vendors; however,
MLPerf focuses on a fixed set of reference models (ResNet, BERT, GPT-3,
Stable Diffusion) and does not cover the full spectrum of LLM training
algorithms.  Several independent studies have compared
ROCm and CUDA on transformer workloads using frameworks such as PyTorch and
JAX, finding convergent results on most standard operations but noting gaps
in operator coverage and compilation overhead.  Our work extends this by
covering \emph{every stage} of the LLM development pipeline rather than a
single training or inference workload.

\paragraph{Framework-level portability.}
PyTorch's \texttt{torch.accelerator} and the \texttt{torch.amp} unified API
represent the framework's answer to vendor portability~\citep{paszke2019pytorch}.
Earlier work on OpenCL-based portability (e.g.\ PlaidML) did not achieve
the performance parity of native vendor paths.  The ROCm port of PyTorch,
maintained by AMD, has reached a level of maturity where the same Python
source compiles and runs correctly on both CUDA and ROCm; our methodology
exploits this via the thin \texttt{DeviceCtx} wrapper.

\paragraph{Distributed training systems.}
Megatron-LM~\citep{narayanan2021efficient} and DeepSpeed~\citep{rajbhandari2020zero} are the
dominant frameworks for large-scale distributed LLM training, with explicit
AMD/ROCm support added in recent releases.  Both frameworks have published
cross-vendor results, but focus on tensor-parallel and pipeline-parallel
regimes (billions of parameters) that are out of scope for MiniMind's
DDP-only, single-node setup.

\paragraph{RLHF and alignment training.}
The PPO-based RLHF paradigm introduced by \citet{ouyang2022training} and the DPO
simplification of \citet{rafailov2023direct} have been implemented in frameworks such as
TRL and OpenRLHF, some of which have begun adding ROCm support.  MiniMind's
hand-written RLHF trainers (PPO, GRPO, DPO) provide an unmediated comparison
that is not filtered through framework-level abstractions.

\paragraph{MoE model training.}
Mixture-of-experts scaling was studied extensively in~\citet{gshard,switch_transformer}.
Vendor-specific MoE kernels (e.g.\ Megablocks~\citep{megablocks}) have been
developed to reduce token-dispatch overhead on CUDA.  ROCm equivalents via
Composable Kernel are an active area of development.  Our study uses
MiniMind's pure-PyTorch MoE dispatch, which avoids custom kernels and
therefore measures the \emph{baseline} cross-vendor MoE performance
achievable without vendor-specific tuning.

\paragraph{Inference engine benchmarks.}
vLLM~\citep{kwon2023efficient} introduced PagedAttention for efficient KV-cache management;
SGLang~\citep{zheng2024sglang} introduced RadixAttention for prefix sharing.  Both
engines maintain ROCm branches.  Benchmarks comparing vLLM and SGLang on
CUDA exist~\citep{zheng2024sglang}, but systematic cross-vendor comparison using
identical converted weights on the same model is, to our knowledge, not
previously published.

\paragraph{Small models as research vehicles.}
The use of small, reproducible models for algorithm research has precedent in
works such as nanoGPT~\citep{nanogpt} and TinyLlama~\citep{tinyllama}.
MiniMind extends this tradition by including the full post-training and RLHF
pipeline as first-class components, making it suitable for cross-vendor
algorithm coverage studies.
